\chapter{事前分布と事後分布の考え方}
    \section{離散型}
        \begin{figure}[H]
            \centering
            \includegraphics[keepaspectratio, scale=0.3]
            {\currfiledir/imgs/prior-and-posterior/graph-discrete.eps}
            \caption{事象$\theta_i$と$X$}
        \end{figure}
        巷で使われている式は次のようにして導かれる。
        \[\Pr{\theta_i|X} = \frac{\Pr{\theta_i\cap X}}{\Pr{X}} = \frac{\Pr{X|\theta_i}\Pr{\theta_i}}{\Pr{X}} = \frac{\Pr{X|\theta_i}\Pr{\theta_i}}{\sum_{i=1}^n{\Pr{X\cap\theta_i}}} = \frac{\Pr{X|\theta_i}\Pr{\theta_i}}{\sum_{i=1}^n{\Pr{X|\theta_i}\Pr{\theta_i}}}\]
    \section{連続型}
        \label{subsec:連続型の事後分布}
        \begin{figure}[H]
            \centering
            \includegraphics[keepaspectratio, scale=0.4]
            {\currfiledir/imgs/prior-and-posterior/graph-continuous.eps}
            \caption{事象$\theta$と$x$の同時確率密度}
        \end{figure}
        巷で使われている式を導出する。
        連続型の場合は少し厄介だ。
        まずいくつか定義を行う。
        \begin{itemize}
            \item $p(x,\theta)$: $x,\theta$の同時確率密度関数。$\integrate{\xmin}{\xmax}{\integrate{\thetamin}{\thetamax}{p(x,\theta)}{}{\theta}}{}{x} = 1$
            \item $\pi(\theta) \coloneqq p_\Theta(\theta) \coloneqq \integrate{\xmin}{\xmax}{p(x,\theta)}{}{x}$ \quad: $\Theta$の周辺確率密度。これが所謂「\textbf{事前分布}」。
            \item $f(x) \coloneqq p_X(x) \coloneqq \integrate{\thetamin}{\thetamax}{p(x,\theta)}{}{\theta}$ \quad: $X$の周辺確率密度。
            \item $f(x|\theta) \coloneqq p_X(x|\Theta=\theta) = \frac{p(x,\theta)}{\pi(\theta)}$ \quad: $X$の条件付き確率密度。これが所謂「\textbf{尤度}」。$\integrate{\xmin}{\xmax}{f(x|\theta)}{}{x} = 1$
            \item $\pi(\theta|x) \coloneqq p_\Theta(\theta|X=x) = \frac{p(x,\theta)}{f(x)}$ \quad: $\Theta$の条件付き確率密度。これが所謂「\textbf{事後分布}」。$\integrate{\thetamin}{\thetamax}{\pi(\theta|x)}{}{\theta} = 1$
        \end{itemize}
        以上の定義の下で、巷で使われている式が導かれる。
        \[\pi(\theta|x) = \frac{p(x,\theta)}{f(x)} = \frac{f(x|\theta)\pi(\theta)}{\integrate{\thetamin}{\thetamax}{p(x,\theta)}{}{\theta}} = \frac{f(x|\theta)\pi(\theta)}{\integrate{\thetamin}{\thetamax}{f(x|\theta)\pi(\theta)}{}{\theta}}\]
        上式から分かるように、尤度$f(x|\theta)$には\underline{定数倍の任意性}があるので、単純に$f(x|\theta)\coloneqq p(x,\theta)$としても上式は成り立つ。
        \par
        実際の計算では、$\pi(\theta)$と$f(x|\theta)$がわかれば、その積を$\theta$に関して規格化すれば$\pi(\theta|x)$が求まる。
        \par
        ベイズ推定では$\pi(\theta)$を我々が\underline{勝手に}決める。
        そして$\Theta=\theta$というパラメータの下で母集団からのサンプルが従う分布$f(x|\theta)$を計算して先述の定義式の4つ目から$p(x,\theta)$を逆算し、最後に事後分布を得ている。